\documentclass[../DoAn.tex]{subfiles}
\begin{document}

Chương 3 đã cung cấp hệ thống nền tảng lý thuyết khoa học nền móng từ không gian vector đặc trưng khuôn mặt, tiêu chuẩn bảo mật tài liệu quốc gia ICAO Doc 9303 cho đến các giải pháp truyền thông phân tán thời gian thực. Tiếp nối các cơ sở lý thuyết vững chắc đó, Chương 4 sẽ chi tiết hóa quy trình thiết kế kiến trúc hệ thống, đặc tả thiết kế chi tiết gói, cấu trúc cơ sở dữ liệu quan hệ, đồng thời tổng hợp các kết quả thực nghiệm xây dựng ứng dụng, kịch bản kiểm thử và mô hình triển khai thực tế của hệ thống điểm danh thông minh tích hợp xác thực kép.

\section{Thiết kế kiến trúc}
\label{section:4.1}

\subsection{Lựa chọn kiến trúc phần mềm}

Hệ thống điểm danh thông minh áp dụng mô hình kiến trúc hướng dịch vụ phân tán kết hợp kiến trúc phân tầng hệ thống tại từng dịch vụ thành phần. Sự kết hợp này nhằm giải quyết trực tiếp các yêu cầu cô lập lỗi, phân tách cổng mạng an ninh và tối ưu hóa năng lực tính toán xử lý đồng thời đã được xác định tại Chương 2. Toàn bộ hệ thống được chia thành năm thành phần độc lập: Phân hệ trạm đầu cuối Android Client, Phân hệ trung tâm điều phối Spring Boot Backend, Phân hệ tính toán sinh trắc học Python Flask AI Service, Phân hệ quản trị trực tuyến React Dashboard, và Phân hệ chấp hành IoT ESP32 Controller.

Áp dụng vào hệ thống thực tế, thành phần trung tâm Spring Boot Backend được tổ chức theo kiến trúc phân tầng chuẩn để quản lý dữ liệu, bao gồm ba lớp chính:
\begin{itemize}
    \item \textbf{Lớp Giao diện lập trình (\texttt{Controller Layer}):} Tiếp nhận các yêu cầu HTTP REST API mã hóa từ ứng dụng di động trạm đầu cuối hoặc kết nối song công WebSocket từ giao diện Web Dashboard, thực hiện kiểm tra định dạng dữ liệu đầu vào.
    \item \textbf{Lớp Nghiệp vụ chính (\texttt{Service Layer}):} Chứa toàn bộ logic xử lý trung tâm của hệ thống bao gồm điều phối luồng xác thực hai lớp sinh trắc học, kiểm tra ma trận phân quyền truy cập và phát thông điệp điều khiển phần cứng. Thành phần này được hiện thực hóa qua các Interface như \texttt{AttendanceService}, \texttt{DeviceService} và các Class triển khai tương ứng như \texttt{AttendanceServiceImpl}, \texttt{DeviceServiceImpl}.
    \item \textbf{Lớp Truy cập dữ liệu (\texttt{Repository Layer}):} Thực thi ánh xạ đối tượng và giao tiếp trực tiếp với cơ sở dữ liệu MySQL, được định nghĩa thông qua các thực thể dữ liệu kế thừa từ Spring Data JPA như \texttt{AttendanceRepository}, \texttt{UserRepository} và \texttt{AccessPermissionRepository}.
\end{itemize}

\subsection{Thiết kế tổng quan}

Biểu đồ gói hệ thống được tổ chức phân rã theo các tầng kiến trúc từ trên xuống dưới, tuân thủ nghiêm ngặt quy tắc kiểm soát phụ thuộc: các gói tầng trên chỉ được phép gọi dịch vụ từ gói tầng dưới trực tiếp, tuyệt đối không phụ thuộc ngược và không phụ thuộc bắc cầu bỏ qua tầng trung gian.

\begin{figure}[tbhp]
    \centering
    \includegraphics[width=0.85\linewidth]{Hinhve/Picture1.png}
    \caption{Biểu đồ phụ thuộc gói tổng quan trong hệ thống Backend trung tâm}
    \label{fig:package_dep}
\end{figure}

Sự phụ thuộc giữa năm gói chức năng chính tại Máy chủ trung tâm được quy định chi tiết:
\begin{itemize}
    \item \textbf{Gói \texttt{com.system.config}:} Tầng cấu hình hệ thống, thiết lập các tham số môi trường, cấu hình bộ lọc an ninh Spring Security, mã hóa bảo mật đường truyền TLS/SSL cho MQTT Broker và WebSocket endpoint. Gói này độc lập và cung cấp hạ tầng chạy cho toàn bộ hệ thống.
    \item \textbf{Gói \texttt{com.system.controller}:} Tầng tiếp nhận request, phụ thuộc trực tiếp vào gói nghiệp vụ \texttt{com.system.service} để đẩy luồng dữ liệu xử lý logic ngầm.
    \item \textbf{Gói \texttt{com.system.service}:} Tầng chứa logic nghiệp vụ lõi, phụ thuộc vào gói truy cập dữ liệu \texttt{com.system.repository} để truy vấn trạng thái và gói \texttt{com.system.dto} để chuẩn hóa cấu trúc dữ liệu payload truyền thông điệp nội bộ.
    \item \textbf{Gói \texttt{com.system.repository}:} Tầng giao tiếp cơ sở dữ liệu, phụ thuộc vào gói \texttt{com.system.entity} chứa các lớp ánh xạ quan hệ thực thể trực tiếp xuống các bảng của MySQL.
\end{itemize}

\subsection{Thiết kế chi tiết gói}

Thiết kế chi tiết cấu trúc mối quan hệ giữa các lớp trong gói nghiệp vụ ghi nhận điểm danh hai lớp thể hiện rõ nét sự phân tách trách nhiệm thông qua các mối quan hệ thực thi (\textit{Implementation}), phụ thuộc (\textit{Dependency}) và kết hợp (\textit{Association}).

\begin{figure}[tbhp]
    \centering
    \includegraphics[width=0.85\linewidth]{Hinhve/Picture2.png}
    \caption{Biểu đồ thiết kế các lớp hạt nhân trong gói dịch vụ điểm danh}
    \label{fig:class_detail_package}
\end{figure}

Sơ đồ thiết kế vạch rõ cơ chế vận hành ngầm của gói: Lớp \texttt{AttendanceController} thực hiện mối quan hệ kết hợp trực tiếp với Interface \texttt{AttendanceService}. Khi nhận request điểm danh, lớp này gọi phương thức xử lý của lớp triển khai \texttt{AttendanceServiceImpl}. Lớp thực thi nghiệp vụ \texttt{AttendanceServiceImpl} đồng thời phụ thuộc vào \texttt{AttendanceRepository} để ghi nhật ký, kết hợp quan hệ phụ thuộc chéo sang \texttt{MqttGateway} để xuất bản lệnh mở cửa xuống thiết bị IoT và \texttt{WebSocketGateway} để phát broadcast cập nhật trực tuyến giao diện người dùng web.

\section{Thiết kế chi tiết}
\label{section:4.2}

\subsection{Thiết kế giao diện}

Hệ thống hướng tới hai môi trường hiển thị đặc thù tương ứng với hai nhóm tác nhân: Môi trường Trạm đầu cuối di động Android Client và Môi trường Bảng quản trị Web Dashboard .

Quy chuẩn thiết kế giao diện được thống nhất nghiêm ngặt trên toàn hệ thống:
\begin{itemize}
    \item \textbf{Phối màu chủ đạo:} Sử dụng phong cách thiết kế giao diện tối cho hệ thống trạm Android và đồ thị Web Dashboard quản trị an ninh nhằm tăng khả năng tập trung. Màu nền chính áp dụng mã màu \texttt{\#0F1117}, cấu phần bề mặt áp dụng \texttt{\#1A1D27}, màu nhấn chức năng an toàn hiển thị màu xanh lá \texttt{\#3ECF8E}, màu cảnh báo vi phạm hiển thị màu đỏ \texttt{\#F76F6F}.
    \item \textbf{Quy chuẩn nút và điều khiển:} Các nút tương tác kích hoạt chức năng hệ thống bắt buộc có chiều cao tối thiểu 44dp trên di động để tối ưu thao tác chạm vật lý. Trạng thái phản hồi thông điệp hệ thống (\textit{Toast/Alert notification}) luôn được neo cố định tại góc trên bên phải màn hình đối với giao diện Web và hiển thị dạng banner tràn viền ngay dưới khung camera đối với ứng dụng Trạm di động Android.
\end{itemize}

\subsection{Thiết kế lớp}

Đoán tiến hành đặc tả cấu trúc phương thức và thuộc tính chi tiết cho lớp dịch vụ hạt nhân điều phối trung tâm hệ thống \texttt{AttendanceServiceImpl}:
\begin{itemize}
    \item \textbf{Các thuộc tính chủ đạo:} 
    \newline \texttt{private final AttendanceRepository attendanceRepo;} (quản lý lưu trữ); 
    \newline \texttt{private final RestTemplate restTemplate;} (gọi dịch vụ mạng đối soát AI); 
    \newline \texttt{private final MqttGateway mqttGateway;} (truyền thông IoT).
    \item \textbf{Phương thức cốt lõi:} 
    \newline \texttt{public AttendanceResponse processAttendance(AttendanceRequest request);}
    \item \textit{Ý nghĩa:} Tiếp nhận payload chứa thông tin định danh và chuỗi hình ảnh selfie. Nếu nhân sự đã kích hoạt (\texttt{ACTIVE}), hệ thống sử dụng ảnh gốc từ kho lưu trữ trung tâm để thực hiện đối soát, cho phép bỏ qua bước đọc NFC vật lý. Chu trình đối soát thực thi qua dịch vụ AI, tính toán khoảng cách toán học so sánh với ngưỡng an ninh $0.55$, cập nhật trạng thái cơ sở dữ liệu và kích hoạt lệnh mở chốt rơ-le vật lý thông qua giao thức MQTT đa kênh (Direct \& Broadcast MQTT). Cơ chế này đảm bảo lệnh điều khiển luôn tới được thiết bị chấp hành ngay cả khi có sự cố định tuyến lẻ.
\end{itemize}

Để minh họa luồng truyền thông điệp chi tiết giữa các thực thể đối tượng phần mềm động trong hệ thống, biểu đồ tuần tự cho Use Case ghi nhận điểm danh hai lớp kết hợp điều khiển thiết bị IoT ngoại vi được thiết lập chi tiết.

Luồng thông điệp diễn ra tuần tự từ Trạm đầu cuối Android, qua lõi bộ điều phối Spring Boot Server, tương tác chéo dịch vụ AI Python Server, ghi đĩa dữ liệu MySQL, và cuối cùng bắn lệnh điều khiển MQTTS giải phóng chốt cửa vật lý trên ESP32 như mô tả tại hình vẽ trên.

\subsection{Thiết kế cơ sở dữ liệu}

Mô hình dữ liệu quan hệ của hệ thống được thiết kế tối ưu hóa cấu trúc chuẩn hóa 3NF, đảm bảo tính toàn vẹn thực thể và toàn vẹn tham chiếu nghiêm ngặt, biểu diễn trực quan qua biểu đồ quan hệ ERD chặt chẽ.

\begin{figure}[tbhp]
    \centering
    \includegraphics[width=0.85\linewidth]{Hinhve/ERD.png}
    \caption{Biểu đồ mô hình quan hệ thực thể cơ sở dữ liệu hệ thống}
    \label{fig:erd_database}
\end{figure}

Kiến trúc cơ sở dữ liệu MySQL bao gồm các bảng vật lý cốt lõi liên kết chặt chẽ:
\begin{itemize}
    \item \textbf{Bảng \texttt{users}:} Lưu trữ tài khoản quản lý web gồm \texttt{id}, \texttt{username}, \texttt{password} (Bcrypt), \texttt{full\_name}, \texttt{email}, \texttt{is\_verified}, \texttt{is\_active}, \texttt{created\_at}. Liên kết với \texttt{user\_roles} để gán vai trò RBAC.
    \item \textbf{Bảng \texttt{employees}:} Dữ liệu nhân sự cốt lõi gồm \texttt{id}, \texttt{cccd} (UNIQUE, 12 ký tự), \texttt{full\_name}, \texttt{birthday}, \texttt{gender}, \texttt{image\_raw\_url} (ảnh DG2 từ chip), \texttt{face\_vector} (JSON format 512D), \texttt{backup\_photo\_embedding} (JSON), \texttt{department}, \texttt{position}, \texttt{email}. Đây là bảng trung tâm chứa toàn bộ dữ liệu nhân sự kèm vector sinh trắc học.
    \item \textbf{Bảng \texttt{attendance\_logs}:} Nhật ký điểm danh gồm \texttt{id}, \texttt{cccd}, \texttt{captured\_name}, \texttt{device\_code}, \texttt{scan\_time}, \texttt{selfie\_image} (Dữ liệu Base64 ảnh selfie), \texttt{score} (Điểm cosine similarity), \texttt{status} (SUCCESS|FAILED). Thêm các cột kỹ thuật theo dõi luồng 2 lớp: \texttt{layer1\_passed}, \texttt{layer1\_score}, \texttt{selfie\_vs\_db\_score}, \texttt{selfie\_vs\_chip\_score}, \texttt{matched}, \texttt{access\_granted}, \texttt{alert\_type} (Cảnh báo Anti-passback) và \texttt{notes} để lưu vết logic tự phục hồi.
    \item \textbf{Bảng \texttt{access\_permissions}:} Ma trận phân quyền gồm \texttt{id}, \texttt{employee\_id} (FK), \texttt{device\_id} (FK), \texttt{start\_date}, \texttt{end\_date}, \texttt{start\_time}, \texttt{end\_time}, \texttt{is\_active}. Định nghĩa constraint UNIQUE(employee\_id, device\_id) để tránh trùng lặp quyền truy cập.
    \item \textbf{Bảng \texttt{devices}:} Thiết bị cửa gồm \texttt{id}, \texttt{device\_code} (UNIQUE, 50 ký tự), \texttt{location\_name}, \texttt{is\_online}, \texttt{last\_sync}.
    \item \textbf{Bảng \texttt{user\_roles}:} Bảng liên kết (N:M) giữa \texttt{users} và \texttt{roles}, định nghĩa PRIMARY KEY(user\_id, role\_id).
\end{itemize}

\section{Xây dựng ứng dụng}
\label{section:4.3}

\subsection{Thư viện và công cụ sử dụng}

Toàn bộ quá trình nghiên cứu, lập trình hiện thực hóa mã nguồn hệ thống phân tán được triển khai đồng bộ dựa trên danh mục các công cụ, môi trường phát triển và thư viện chuyên dụng được tổng hợp chi tiết tại Bảng \ref{table:tools_development}.

\begin{table}[tbhp]
\centering
\small
\renewcommand{\arraystretch}{1.3}
\caption{Danh sách môi trường, thư viện và công cụ tích hợp phát triển ứng dụng}
\label{table:tools_development}
\begin{tabular}{|p{3.2cm}|p{4.2cm}|p{6.8cm}|}
\hline
\textbf{Mục đích sử dụng} & \textbf{Tên công cụ / Thư viện} & \textbf{Phiên bản / Địa chỉ URL hệ thống} \\
\hline
Môi trường lập trình lõi & IntelliJ IDEA Ultimate & v2023.3 — https://www.jetbrains.com/idea/ \\
\hline
Framework Backend trung tâm & Java Spring Boot & v3.2.3 — https://spring.io/projects/spring-boot \\
\hline
Dịch vụ xử lý sinh trắc học & Python Flask AI Server & v3.0.2 — https://flask.palletsprojects.com/ \\
\hline
Thư viện nhận diện học sâu & InsightFace (ArcFace Engine) & v0.7.3 — https://github.com/deepinsight/insightface \\
\hline
Môi trường nhúng thiết bị & FaceNet-PyTorch (ONNX) & v2.5.3 — https://github.com/timesler/facenet-pytorch \\
\hline
Framework Trạm đầu cuối & Google Flutter (Dart v3.2) & v3.19.2 — https://flutter.dev/ \\
\hline
Giao tiếp trường gần NFC & Thư viện \texttt{dmrtd} Core & v0.4.1 — https://pub.dev/packages/dmrtd \\
\hline
Thư viện giao diện Web & React JS Ecosystem & v18.2.0 — https://react.dev/ \\
\hline
Hạ tầng môi giới mạng IoT & HiveMQ Cloud (SSL/TLS) & Cổng bảo mật 8883 — https://www.hivemq.com/ \\
\hline
Hệ quản trị cơ sở dữ liệu & MySQL Community Server & v8.0.36 — https://www.mysql.com/ \\
\hline
\end{tabular}
\end{table}

\subsection{Kết quả đạt được}

Sản phẩm hoàn thiện đóng gói của đồ án tốt nghiệp bao gồm ba cấu phần phần mềm và một cấu phần phần cứng nhúng độc lập: Tệp tin mã máy đóng gói ứng dụng Android (\texttt{app-release.apk}) cài đặt trực tiếp tại máy trạm đầu cuối; Thư mục mã nguồn thực thi Backend trung tâm Java Spring Boot đóng gói dưới dạng tệp tin lưu trữ tiêu chuẩn \texttt{JAR}; Phân hệ dịch vụ trí tuệ nhân tạo Python được ảo hóa cô lập qua container \texttt{Docker}; và tệp tin mã nhúng nhị phân \texttt{firmware.bin} nạp trực tiếp xuống bo mạch vi điều khiển ESP32.

Hệ thống tiến hành thống kê định lượng chi tiết về quy mô phân bố mã nguồn vật lý được tổng hợp cụ thể trong Bảng \ref{table:source_stats}.

\begin{table}[tbhp]
\centering
\small
\renewcommand{\arraystretch}{1.3}
\caption{Bảng thông số thống kê định lượng quy mô mã nguồn ứng dụng}
\label{table:source_stats}
\begin{tabular}{|p{4cm}|c|c|p{4.5cm}|}
\hline
\textbf{Phân hệ ứng dụng} & \textbf{Số dòng code} & \textbf{Số lượng lớp} & \textbf{Dung lượng sản phẩm đóng gói} \\
\hline
Spring Boot Backend & ~4,500 lines & 32 Classes & 42 MB (File \texttt{JAR} nhúng Tomcat) \\
\hline
Flutter Android Station & ~6,200 lines & 45 Classes & 18.5 MB (File APK phân phối) \\
\hline
Python Flask AI Service & ~1,800 lines & 12 Modules & 1.2 GB (Gồm trọng số Weights mạng) \\
\hline
React Web Dashboard & ~3,100 lines & 24 Components & 3.8 MB (Thư mục Production Build) \\
\hline
ESP32 Firmware C++ & ~850 lines & 4 Modules & 680 KB (Tệp tin nhị phân \texttt{bin}) \\
\hline
\end{tabular}
\end{table}

\subsection{Minh họa các chức năng chính}

Hệ thống đã triển khai xây dựng thành công và vận hành ổn định các chức năng hạt nhân trung tâm bao gồm:
\begin{itemize}
    \item \textbf{Chức năng Ghi nhận điểm danh hai lớp tại Trạm Android:} Giao diện màn hình trạm luôn mở sẵn bộ lọc camera giám sát thời gian thực. Khi đặt thẻ CCCD vào vùng cảm biến sau lưng thiết bị, ứng dụng kích hoạt tiến trình bóc tách tệp phân vùng DG2 ngầm để giải mã chân dung gốc. Đồng thời giao diện hiển thị hiệu ứng quét sóng ngắn hình tròn màu xanh lục quanh khuôn mặt nhân sự, tự động chụp selfie và thực hiện thuật toán tính toán khoảng cách cosine lớp thứ nhất ngay tại trạm đầu cuối.
    \item \textbf{Chức năng Giám sát dòng người ra vào trực tuyến trên Web Dashboard:} Giao diện quản trị viên React sử dụng mô hình bảng biểu động kết hợp thư viện kết nối thời gian thực WebSocket. Ngay khi một lượt điểm danh được phê duyệt thành công tại server, một dòng bản ghi mới lập tức được đẩy tự động chèn vào đầu danh sách hiển thị với đầy đủ thông tin so khớp ảnh gốc từ chip của Bộ Công an đặt song song cạnh ảnh chụp selfie thực tế tại hiện trường. Đặc biệt, hệ thống tích hợp cơ chế nhận diện chiều di chuyển trực quan thông qua bộ mã màu chuẩn: màu xanh biểu tượng cho luồng "Vào (IN)" và màu đỏ cho luồng "Ra (OUT)", giúp người quản lý dễ dàng bao quát trạng thái an ninh của từng cửa kiểm soát mà không cần đọc chi tiết văn bản. Hệ thống cũng cung cấp bảng điều khiển cấu hình chiều di chuyển Admin-only, bảo đảm tính bảo mật trong việc thiết lập hạ tầng thực địa.
\end{itemize}

\section{Kiểm thử}
\label{section:4.4}

Đồ án áp dụng kết hợp hai phương pháp luận kiểm thử phần mềm khoa học bao gồm: Kiểm thử hộp đen (\textit{Black-box Testing}) áp dụng tại tầng giao diện để đánh giá luồng nghiệp vụ đầu cuối, và Kiểm thử tích hợp (\textit{Integration Testing}) đánh giá năng lực tương tác giao thức mạng liên nền tảng. Tiến trình thực nghiệm thiết kế ba trường hợp kiểm thử (Test Cases) hạt nhân đại diện cho các điều kiện biên của hệ thống.

\subsection{Đặc tả kịch bản kiểm thử hệ thống}

\textbf{Test Case 1: Thực hiện điểm danh hai lớp thành công với thẻ chính chủ}
\begin{itemize}
    \item \textit{Điều kiện đầu vào:} Nhân sự sử dụng thẻ CCCD hợp lệ còn hiệu lực, tài khoản trạng thái \texttt{ACTIVE}, đứng trước camera trạm Android trong khung giờ được phân quyền truy cập.
    \textit{Các bước thực hiện:} (1) Chạm thẻ vào vùng cảm biến NFC của thiết bị trạm; (2) Nhìn thẳng vào camera chụp ảnh selfie thực tế; (3) Hệ thống tự động thực thi luồng đối soát hai tầng ngầm và truyền tải thông điệp.
    \item \textit{Kết quả kỳ vọng:} Ứng dụng trạm hiển thị thông báo thành công màu xanh lá, Web Dashboard lập tức chèn dòng nhật ký điểm danh mới thời gian thực, mạch ESP32 nhận thông điệp MQTT (qua topic định danh hoặc topic broadcast dự phòng) kích rơ-le giải phóng chốt cửa vật lý mở ra trong 3 giây.
    \item \textit{Kết quả thực tế:} Đạt tiêu chuẩn cấu hình kỳ vọng (Pass). Hệ thống ghi nhận độ trễ xử lý đầu cuối tổng thể từ lúc chụp ảnh đến lúc mở khóa dao động ổn định trong khoảng 2.8 -- 3.8 giây tùy điều kiện mạng.
\end{itemize}

\textbf{Test Case 2: Từ chối quyền truy cập do giả mạo sinh trắc học khuôn mặt}
\begin{itemize}
    \item \textit{Điều kiện đầu vào:} Người thực hiện điểm danh sử dụng thẻ CCCD của một nhân sự khác (thẻ thật, chip thật) nhưng cố tình sử dụng một tấm ảnh in chân dung hoặc thiết bị video phát lại khuôn mặt của chủ thẻ trước camera trạm.
    \textit{Các bước thực hiện:} (1) Chạm thẻ CCCD hợp lệ vào vùng cảm biến; (2) Đưa ảnh in hoặc video khuôn mặt chủ thẻ đối diện camera để thực hiện chụp selfie đối soát.
    \item \textit{Kết quả kỳ vọng:} Hệ thống tính toán điểm số tương đồng hình học không gian vector thực tế nằm dưới ngưỡng cấu hình an ninh được tuning sẵn. Trạm Android phát âm thanh cảnh báo lỗi xác thực, ghi nhận log trạng thái \texttt{FAILED} về máy chủ trung tâm, tuyệt đối giữ chốt rơ-le vật lý ESP32 khóa chặt cửa.
    \item \textit{Kết quả thực tế:} Đạt tiêu chuẩn cấu hình kỳ vọng (Pass). Mô hình ArcFace tại máy chủ dịch vụ AI bóc tách chính xác sai biệt hình học và chặn đứng hành vi gian lận danh tính.
\end{itemize}

\textbf{Test Case 3: Cơ chế xử lý dự phòng ngoại tuyến khi mất kết nối Internet}
\begin{itemize}
    \item \textit{Điều kiện đầu vào:} Thiết bị trạm đầu cuối Android bị ngắt kết nối mạng Wifi/4G cục bộ (mất liên lạc hoàn toàn với Spring Boot Server trung tâm).
    \textit{Các bước thực hiện:} Nhân viên chạm thẻ CCCD hợp lệ và chụp ảnh selfie thực hiện thao tác điểm danh tại trạm đầu cuối trong điều kiện ngoại tuyến.
    \item \textit{Kết quả kỳ vọng:} Ứng dụng Android không crash hệ thống, tự động kích hoạt bộ lọc chạy đối soát sinh trắc học lớp thứ nhất tại chỗ (Local MobileFaceNet). Nếu hợp lệ, ứng dụng đóng gói payload mã hóa và tự động đẩy bản ghi vào hàng đợi cơ sở dữ liệu SQLite cục bộ, hiển thị trạng thái chờ đồng bộ. Khi khôi phục kết nối mạng, hệ thống tự động kích hoạt luồng đẩy hàng loạt bản ghi về máy chủ trung tâm.
    \item \textit{Kết quả thực tế:} Đạt tiêu chuẩn cấu hình kỳ vọng (Pass). Toàn bộ dữ liệu nhật ký được bảo toàn tính toàn vẹn tuyệt đối, không xảy ra hiện tượng thất thoát dữ liệu chấm công của tổ chức.
\end{itemize}

\textbf{Test Case 4: Kiểm soát Anti-passback và quy trình ra vào đúng trình tự}
\begin{itemize}
    \item \textit{Điều kiện đầu vào:} Nhân sự đã thực hiện điểm danh Vào (IN) thành công, hoặc nhân sự cố tình quét điểm danh Ra (OUT) khi chưa có lịch sử vào trong ngày.
    \textit{Các bước thực hiện:} (1) Nhân viên quét thẻ tại trạm OUT khi trước đó chưa quét tại trạm IN; (2) Nhân viên quét thẻ Ra liên tiếp hai lần tại các trạm OUT.
    \item \textit{Kết quả kỳ vọng:} Hệ thống truy vấn lịch sử tức thời, phát hiện vi phạm trình tự logic. Trạm đầu cuối hiển thị thông báo "🚫 Anti-Passback Alert", trạng thái bản ghi ghi nhận \texttt{FAILED} với lý do cụ thể (Chưa quét vào hoặc Quét ra liên tiếp), và tuyệt đối không kích hoạt mở cửa.
    \item \textit{Kết quả thực tế:} Đạt tiêu chuẩn cấu hình kỳ vọng (Pass). Hệ thống ngăn chặn thành công 100\% các hành vi đi ngược luồng hoặc gian lận quy trình ra vào.
\end{itemize}

\subsection{Tổng kết thực nghiệm kiểm thử}

Hệ thống đã triển khai thực thi tổng cộng 52 kịch bản trường hợp kiểm thử hộp đen và kiểm thử tích hợp API (bao gồm cả các kịch bản kiểm tra logic Anti-passback mới). Kết quả thực nghiệm ghi nhận 50 kịch bản đạt tiêu chuẩn (đạt tỉ lệ thành công 96.15\%), 2 kịch bản không đạt liên quan đến lỗi Timeout đường truyền mạng truyền tải thông điệp khi bo mạch ESP32 bị mất nguồn cấp vật lý đột ngột. Hệ thống xử lý lỗi ngoại lệ bằng cách tự động cấu hình cơ chế lắng nghe phản hồi ngược và phát cảnh báo âm thanh lên Web Dashboard cho nhân viên an ninh giám sát thực địa.

\section{Triển khai thử nghiệm và đánh giá hệ thống}
\label{section:4.5}

Hệ thống điểm danh thông minh tích hợp xác thực hai lớp được triển khai thử nghiệm thực tế trong môi trường hạ tầng phân tán mạng nội bộ kết hợp điện toán đám mây. Cấu hình chi tiết các phân hệ máy chủ vận hành hệ thống được thiết lập khắt khe: Máy chủ trung tâm điều phối Spring Boot lõi và Cơ sở dữ liệu MySQL 8.0 được triển khai trực tiếp trên hạ tầng phần cứng Dell PowerEdge (Cấu hình CPU Intel Xeon 8 Cores, dung lượng bộ nhớ RAM 16GB, hệ điều hành Ubuntu Server v22.04 LTS); Phân hệ dịch vụ trí tuệ nhân tạo Python Flask AI Server được cấu hình tối ưu năng lực tính toán song song thông qua việc chạy trên máy chủ chuyên dụng hỗ trợ card tăng tốc đồ họa phần cứng GPU NVIDIA Tesla T4 16GB VRAM nhằm tối ưu hóa tốc độ xử lý các toán tử ma trận của mô hình học sâu mạng nơ-ron tích chập.

Kết quả vận hành kiểm toán thực nghiệm của hệ thống sau chu kỳ 4 tuần triển khai chạy thử nghiệm liên tục tại văn phòng cơ quan đạt được các chỉ số định lượng khoa học độ tin cậy cao:
\begin{itemize}
    \item \textbf{Tổng lưu lượng và năng lực chịu tải:} Hệ thống phục vụ tần suất vận hành điểm danh ra vào liên tục của hơn 150 nhân sự chính thức với tổng số lượt tương tác giao tiếp API ghi nhận đạt hơn 600 lượt truy cập điểm danh trung bình mỗi ngày. Hệ thống duy trì hoạt động liên tục 24/7 với chỉ số độ sẵn sàng hệ thống đạt tỷ lệ lý tưởng $99.82\%$.
    \item \textbf{Tốc độ phản hồi và Hiệu năng tính toán:} Chu kỳ thời gian xử lý trung bình cho một lượt điểm danh hai lớp khép kín ghi nhận thực tế đạt mức \textbf{2.8 - 3.8 giây} (trong đó tiến trình bắt tay bảo mật NFC đọc chip CCCD tiêu tốn khoảng 1.5 - 2.0 giây, chụp ảnh selfie 0.5 giây, truyền tải payload và thuật toán trích xuất đối sánh vector đặc trưng tại máy chủ AI tiêu tốn khoảng 0.8 - 1.2 giây). Tốc độ phản hồi hoàn toàn đáp ứng xuất sắc yêu cầu phi chức năng $<5$ giây đặt ra tại Chương 2.
    \item \textbf{Độ chính xác mô hình nhận diện sinh trắc học:} Qua các đợt kiểm thử thực tế đối soát chéo hình ảnh chân dung giữa ảnh camera selfie và ảnh DG2 từ chip CCCD, hệ thống đo lường chỉ số nhận diện ổn định với cấu hình ngưỡng khoảng cách hình học Cosine Similarity tối ưu đạt $0.55$. Kết quả cho thấy mô hình ArcFace có khả năng phân tách biên đặc trưng tốt ngay cả với ảnh chip nén dung lượng thấp.
\end{itemize}

Chỉ số thực nghiệm định lượng khoa học trên chứng minh tính khả thi vượt trội và năng lực ứng dụng thực tế của mô hình kiến trúc điểm danh hai lớp tự trị không lưu kho ảnh gốc dữ liệu tập trung, giải quyết trọn vẹn cả hai yêu cầu mâu thuẫn trong các hệ thống an ninh truyền thống: tối đa hóa mức độ bảo mật chống giả mạo danh tính đồng thời bảo vệ tuyệt đối quyền riêng tư dữ liệu cá nhân của người sử dụng.

\end{document}